%-*-latex-*-
%-*-latex-*-
\newcommand\COURSE{ciss245}
\newcommand\ASSESSMENT{q6501}
\newcommand\ASSESSMENTTYPE{Quiz}
\newcommand\POINTS{\textwhite{xxx/xxx}}

\input{myquizpreamble}
\input{yliow}
\input{\COURSE}
%-*-latex-*-

%-*-latex-*-

%-*-latex-*-
\renewcommand\TITLE{\ASSESSMENTTYPE \ \ASSESSMENT}

\renewcommand\EMAIL{}
\renewcommand\AUTHOR{eagraham1@cougars.ccis.edu}

\textwidth=6in
\begin{document}
\topmatterthree

%------------------------------------------------------------------------------
\nextq
A matrix is a 2D array of integers (for this question).
Here's how you use it:
\begin{Verbatim}[frame=single,fontsize=\footnotesize]
Mat m(2, 3); // m models a 2-by-3 array of integers, all initialized to 0
std::cout << m << '\n'; // prints like this:
                        // 0 0 0
                        // 0 0 0
m(0, 0) = 42;           // sets row 0, column 0 entry to 42
std::cout << m << '\n'; // prints like this:
                        // 42 0 0
                        // 0 0 0
Mat n(2, 2);            // n is 2-by-2
n = m;                  // n is 2-by-3 (like m) and values of m are copied over
std::cout << n << '\n'; // prints like this:
                        // 42 0 0
                        // 0 0 0
Mat p(m);
std::cout << p << '\n'; // prints like this:
                        // 42 0 0
                        // 0 0 0
\end{Verbatim}
Complete the following header file.
Do not implement any methods/functions in the header file
(that's done in the question below).
\\
\ANSWER
\begin{answercode}
// file: Mat.h
#include <iostream>
class Mat
{
public:
Mat(int r, int c);

Mat(const Mat& other);
Mat & operator=(const Mat & other);

~Mat();

int & operator()(int i, int j);
const int & operator()(int i, int j) const;
private:
    int num_rows_; // number of rows
    int num_cols_; // number of columns
    int * p_;      // there's only one reason for this!
};
std::ostream & operator<<(std::ostream &, const Mat &);
\end{answercode}

%------------------------------------------------------------------------------
\nextq
Refer to the question above.
Complete the following implementation file for Matrix class.
Obviously there must be no memory issues.
\\
\ANSWER
\begin{answercode}
// file: Mat.cpp
#include <iostream>
#include "Mat.h"

Mat::Mat(int r, int c)
: num_rows_(r), num_cols_(c), p_(new int[r * c])
{
  for (int i = 0; i < r * c; ++i)
  {
    p_[i] = 0;
  }
}

Mat::Mat(const Mat & other)
: num_rows_(other.num_rows_), num_cols_(other.num_cols_),
p_(new int[other.num_rows_ * other.num_cols_])
{
  for (int i = 0; i < num_rows_ * num_cols_; ++i) {
    p_[i] = other.p_[i];
  }
}

Mat & Mat::operator=(const Mat & other)
{
  if (this != &other)
  {
    delete[] p_;

    num_rows_ = other.num_rows_;
    num_cols_ = other.num_cols_;
    p_ = new int[num_rows_ * num_cols_];

    for (int i = 0; i < num_rows_ * num_cols_; ++i)
    {
      p_[i] = other.p_[i];
    }
  }
  return *this;
}

Mat::~Mat()
{
  delete[] p_;
}

int & Mat::operator()(int i, int j)
{
  return p_[i * num_cols_ + j];
}

const int & Mat::operator()(int i, int j) const
{
  return p_[i * num_cols_ + j];
}

std::ostream & operator<<(std::ostream & cout, const Mat & mat)
{
  for (int i = 0; i < mat.num_rows_; i++)
  {
    for (int j = 0; j < mat.num_cols_; j++)
    {
      cout << mat(i, j);
      if (j < mat.num_cols_ - 1)
      cout << ' ';
    }
    cout << '\n';
  }
  return cout;
}

\end{answercode}
You are advised to use google's asan tool.

%------------------------------------------------------------------------------
%-*-latex-*-
\newpage
\input{instructions}
\end{document}

